\begin{abstract}

Video Instance Segmentation (VIS) faces significant annotation challenges due to its dual requirements of pixel-level masks and temporal consistency labels. While recent unsupervised methods like VideoCutLER eliminate optical flow dependencies through synthetic data, they remain constrained by the synthetic-to-real domain gap. We present AutoQ-VIS, a novel unsupervised framework that bridges this gap through quality-guided self-training. Our approach establishes a closed-loop system between pseudo-label generation and automatic quality assessment, enabling progressive adaptation from synthetic to real videos. Experiments demonstrate state-of-the-art performance with 52.6 $\text{AP}_{50}$ on YouTubeVIS-2019 \texttt{val} set, surpassing the previous state-of-the-art VideoCutLER by 4.4$\%$, while requiring no human annotations. This demonstrates the viability of quality-aware self-training for unsupervised VIS.
%
% We will release the code and models upon acceptance.
The source code of our method is available at \href{https://github.com/wcbup/AutoQ-VIS}{here}.

\end{abstract}
